\centerline{\bf \Huge Текстовые задачи}

\problem{\textbf{1}} Смешали $10 \%$-ный и $25 \%$-ный растворы соли и получили 3 кг 20\%-ного раствора. Какое количество каждого раствора в килограммах было использовано?

\problem{\textbf{2}} В двух сосудах имеется вода разной температуры. Из этой воды составляют смеси. Если отношение объемов воды, взятой из первого и второго сосудов, равно $1: 2$, то температура смеси будет $35^{\circ} \mathrm{C}$, а если $3: 4$, то температура смеси будет $33^{\circ}$ С. Найдите температуру воды в каждом сосуде (считая, что плотность и удельная теплоемкость воды не зависят от температуры).

\problem{\textbf{3}} От пристани $A$ отправились одновременно вниз по течению реки катер и плот. Катер спустился вниз по течению на 96 км, затем повернул обратно и вернулся в $A$ через 14 ч. Найдите скорость катера в стоячей воде и скорость течения реки, если известно, что катер встретил плот на обратном пути на расстоянии 24 км от $A$.

\problem{\textbf{4}} Два поезда отправляются навстречу друг другу из городов $A$ и $B$. Если поезд из города $A$ отправится на 1,5 ч раньше, чем поезд из города $B$, то они встретятся на середине пути. Если оба поезда выйдут одновременно, то через 6 ч они еще не встретятся, а расстояние между ними составит десятую часть первоначального. За сколько часов может проехать каждый поезд расстояние между $A$ и $B$ ?

\problem{\textbf{5}} Велосипедист отправляется с некоторой скоростью из пункта $A$ в пункт $B$, отстоящий от $A$ на расстоянии 60 км. Прибыв в $B$, он сразу же выезжает обратно с той же скоростью, но через 1 ч после выезда из $B$ делает остановку на 20 мин, после чего продолжает путь, увеличив скорость на 4 км/ч. В каких пределах заключена скорость велосипедиста, если известно, что на обратный путь от $B$ до $A$ он затратил времени не более, чем на путь от $A$ до $B$.

\problem{\textbf{6}} Турист отправляется в поход из $A$ в $B$ и обратно и проходит весь путь за 3 ч 41 мин. Дорога из $A$ в $B$ идет сначала в гору, потом по ровному месту, а потом под гору. На каком протяжении дорога тянется по ровному месту, если скорость ходьбы туриста составляет: в гору - 4 км/ч, по ровному месту - 5 км/ч, под гору - 6 км/ч, а расстояние между $A$ и $B$ (по дороге) равно 9 км?


\problem{\textbf{7}} Два тела движутся равномерно по окружности в одну сторону. Первое тело проходит окружность на 3 с быстрее второго и догоняет второе тело каждые полторы минуты. За какое время каждое тело проходит окружность?

\newpage


\problem{\textbf{8}} Двум землекопам было поручено вырыть канаву за 3 ч 36 мин. Однако первый приступил к работе тогда, когда второй уже вырыл треть канавы и перестал копать. В результате канава была вырыта за 8 ч. За сколько часов каждый землекоп может вырыть канаву?

\problem{\textbf{9}} Пять человек выполняют некоторую работу. Первый, второй и третий, работая вместе, выполняют всю работу за 7,5 ч, первый, третий и пятый - за 5 ч, первый, третий и четвертый - за 6 ч, четвертый, второй и пятый - за 4 ч. За какой промежуток времени выполняют эту работу все пять человек вместе?

\problem{\textbf{10}} В бассейн проведены две трубы - подающая и отводящая, причем через первую бассейн наполняется на 2 ч дольше, чем через вторую опорожняется. При заполненном на $1 / 3$ бассейне были открыты обе трубы, и бассейн оказался пустым через 8 ч. За сколько часов, действуя отдельно, первая труба наполняет, а вторая опорожняет бассейн?

\problem{\textbf{11}} В арифметической прогрессии, содержащей девять членов, первый член равен 1 , а сумма всех членов равна 369. Геометрическая прогрессия также имеет девять членов, причем первый и последний ее члены совпадают с соответствующими членами данной арифметической прогрессии. Найдите пятый член геометрической прогрессии.

\problem{\textbf{12}} Сумма первых тринадцати членов арифметической прогрессии равна 130. Известно, что четвертый, десятый и седьмой члены этой прогрессии, взятые в указанном порядке, представляют собой три последовательных члена геометрической прогрессии. Найдите первый член арифметической прогрессии.

\problem{\textbf{13}} Сумма первых десяти членов арифметической прогрессии равна 155, а сумма первых двух членов геометрической прогрессии равна 9 . Найдите эти прогрессии, если первый член арифметической прогрессии равен знаменателю геометрической прогрессии, а первый член геометрической прогрессии равен разности арифметической прогрессии.

\problem{\textbf{14}} Даны две арифметические прогрессии $a_1, a_2, a_3$ и $b_1, b_2, b_3$. Известно, что $a_1+a_2+a_3=b_1+b_2+b_3$, а числа $a_1+b_1$, $a_2+b_2, a_3+b_3$ образуют геометрическую прогрессию. Докажите, что $a_1=b_3, a_2=b_2, a_3=b_1$.

\problem{\textbf{15}} Длины трех отрезков составляют геометрическую прогрессию. При каких значениях знаменателя прогрессии из этих отрезков можно составить треугольник?\\

\centerline{\textbf{225 минут}}